% Native continuation of 260, 288, 315, 323--325 and the summed budget in 330.
\section{Endpoint-tested marked roof operators on the actual collision chain}
\label{sec:wo-endpoint-tests}

We pass from scalar label tests to smooth observations of the physical
collision states.  The estimate below is a weak operator estimate with
specified test norms.  Its proof uses the actual two-endpoint stationary
action, including the full inverse-gradient moment; a derivative of the
expanding forward collision map is never replaced by a bounded derivative.
Throughout, $R\in[R_-,R_+]$, $m\ge1$, and $w=1$ or the original
$W^\varepsilon$ with $\varepsilon>0$ fixed.  The clock, section, labels,
and endpoint memberships are those of \eqref{eq:v76-event}.

Put $\mathcal X=(\mathbb R/2\pi\mathbb Z)\times[-1,1]$.  For a complex
$C^1$ function $a$ on this compact cylinder, including continuous
one-sided derivatives at its boundary, set
\begin{equation}\label{eq:wo-test-norm}
 a_0=\|a\|_\infty,\qquad
 a_1=\max\{\|\partial_\alpha a\|_\infty,
                         \|\partial_p a\|_\infty\},\qquad
 \|a\|_{C^1}=a_0+a_1.
\end{equation}
Use the same convention for $b$.  Let
$\varphi_0,\ldots,\varphi_m$ be real $C^1(\mathcal X)$ functions and
write, on the actual regular $m$-collision chain,
\begin{equation}\label{eq:wo-phase}
 \Psi_m(x)=\sum_{i=0}^m\varphi_i(T_R^ix),\qquad
 l_i=\max\{\|\partial_\alpha\varphi_i\|_\infty,
                    \|\partial_p\varphi_i\|_\infty\},\qquad
 \mathcal L_m=\sum_{i=0}^m l_i(\sigma^i+\sigma^{m-i}).
\end{equation}
These functions are observations of the real orbit.  They impose no
independence assumption and remove no collisions.  Define constants
using precisely the inherited physical Hessian and Jacobi bounds:
\begin{equation}\label{eq:wo-constants}
 C_{\rm test,end}=\Lambda_\circ(\lambda_\circ^{-1}+R_-^{-1}),\qquad
 C_{\rm test,phase}=\frac{\Lambda_\circ}{\lambda_\circ}(J_*+C_c),\qquad
 J_{\rm test}=J_\times/c_*.
\end{equation}
Here $J_*$ and $C_c$ are defined in \eqref{eq:el-constants} and
\eqref{eq:el-observable-constants}, and $J_\times$ in
\eqref{eq:sm-constants}.  The constants contain exactly one original
conditioning factor.  They have no dependence on $m$, a positive
incidence cutoff, a label, a roof value, or a new observation event.

\begin{lemma}[Actual Gauss derivatives of the endpoint tests and orbit phase]
\label{lem:wo-derivatives}
On every regular fixed-Gauss-angle segment, let $t=F(q)$,
$q=(\alpha_0,\alpha_m)$, $g=\nabla_qF$, $z=|g|>0$,
$e=g/z$ and $H=e^{\mathsf T}G^{-1}e$.  Then
\begin{equation}\label{eq:wo-gauss}
 q_t=\frac{G^{-1}e}{zH},\qquad
 z_t=\frac1{zH},\qquad
 \|q_t\|_\infty\le\frac{\Lambda_\circ}{\lambda_\circ z}.
\end{equation}
With the collision states and momenta of this segment,
\begin{align}
 z\bigl|\partial_t a(\alpha_0,p_0)\bigr|
       &\le C_{\rm test,end}a_1,
 &z\bigl|\partial_t b(\alpha_m,p_m)\bigr|
       &\le C_{\rm test,end}b_1,                        \label{eq:wo-endpoint-derivative}\\
 z|\partial_t\Psi_m|&\le C_{\rm test,phase}\mathcal L_m.
                                                        \label{eq:wo-phase-derivative}
\end{align}
In particular, if $l_i\le L$ for every $i$, then
$\mathcal L_m\le2L/(1-\sigma)$, uniformly in the entire collision
count.  There is no inverse incidence in the interior observable
derivatives in this assertion.
\end{lemma}
\begin{proof}
At fixed $e$, differentiation of $\nabla F(q)=ze$ gives
$dq/dz=G^{-1}e$.  Differentiation of the physical action therefore
gives $dt/dz=g\cdot(dq/dz)=zH$.  This proves both identities in
\eqref{eq:wo-gauss}; in particular the factor $H$ must not be
omitted from $q_t$.  The Hessian bounds
\eqref{eq:v109-hessian-bounds} imply
$\Lambda_\circ^{-1}\le H\le\lambda_\circ^{-1}$ and
$\|G^{-1}e\|_\infty\le\lambda_\circ^{-1}$, proving the displayed
inequality.  By the exact first variation
$g=R(-p_0,p_m)$, we have $p_0=-ze_1/R$ and $p_m=ze_2/R$.
Consequently each endpoint momentum derivative is at most
$\Lambda_\circ/(R_-z)$ in absolute value.  The chain rule and the
two actual partial-derivative norms in \eqref{eq:wo-test-norm}
give \eqref{eq:wo-endpoint-derivative}.

For completeness, the momentum derivative at an interior collision
is obtained geometrically, without differentiating a square root
of the incidence.  Let $u_i=n'(\alpha_i)$ be the unit tangent and
$v_i=(x_{i+1}-x_i)/\tau_i$ the unit outgoing flight.  For $i<m$,
the true tangential momentum is $p_i=v_i\cdot u_i$; at $i=m$
use $p_m=v_{m-1}\cdot u_m$, since specular reflection preserves
the tangential component.  For an endpoint direction $\eta$,
\[
 |D_qp_i[\eta]|
 \le \|D_qv_i[\eta]\|_2+|D_q\alpha_i[\eta]|
 \quad(i<m),
\]
with $v_{m-1}$ in the terminal formula.  The unweighted Jacobi
localization \eqref{eq:el-unweighted}, the flight direction bound
\eqref{eq:el-direction}, and
$s_{i+1},s_{i-1}\le\sigma^{-1}s_i$ wherever the indicated
indices exist yield
\begin{equation}\label{eq:wo-interior-momentum}
 |D_qp_i[\eta]|\le C_c s_i\|\eta\|_\infty,
 \qquad |D_q\alpha_i[\eta]|\le J_*s_i\|\eta\|_\infty
 \quad(0\le i\le m).
\end{equation}
This proof also includes $m=1$.  Every flight denominator is the
true length, bounded below by $\tau_-$; no denominator $c_i$
has been introduced.  Applying the chain rule to each
$\varphi_i(\alpha_i,p_i)$ and then inserting $\eta=q_t$ gives
\eqref{eq:wo-phase-derivative}.  Finally
\[
 \sum_{i=0}^m s_i
 =\frac{2(1-\sigma^{m+1})}{1-\sigma}
 \le\frac2{1-\sigma}
\]
proves the uniform phase assertion.  Every differentiation stays
inside the same actual stationary physical branch.  It is not a
bound on $D T_R^i$ with initial momentum held fixed.
\end{proof}

For $\ell=(n,k)$ define the complex physical roof measure by
\begin{equation}\label{eq:wo-measures}
 f^{w;a,b,\Psi}_{\ell,m,R}(t)\,dt
 =(L_{m,R})_*\left(
 w\,\eta_R(x)\eta_R(T_R^mx)
 \one_{\{A_{m,R}=n,\,K^{\rm c}_{m,R}=k\}}
 a(x)b(T_R^mx)e^{i\Psi_m(x)}\,\frac{d\nu(x)}{c_*}
 \right).
\end{equation}
The existence of an $L^1$ density follows from the same exact
physical coarea formula in Theorem~\ref{thm:rv-exact-bv}, by domination with the complete
capped source; the zero-measure critical bulk is exhausted as in Section~\ref{sec:rbv-complete}.
Each density is supported in the unchanged physical interval
$[m\tau_-,m\tau_+]$.

\begin{theorem}[Summed variation with specified smooth trajectory tests]
\label{thm:wo-bv}
For the complete source and all its exact labels,
\begin{align}
 \sum_\ell\|f^{w;a,b,\Psi}_{\ell,m,R}\|_1&\le a_0b_0,
                                                        \label{eq:wo-mass}\\
 \sum_\ell\|D f^{w;a,b,\Psi}_{\ell,m,R}\|_{\rm TV}
 &\le D_{\rm all}^w a_0b_0
 +J_{\rm test}\left[
 C_{\rm test,end}(a_1b_0+a_0b_1)
       +C_{\rm test,phase}\mathcal L_m a_0b_0\right]
 =:\mathcal V_m^w(a,b,\Psi).                         \label{eq:wo-variation}
\end{align}
The label sum is taken before every supremum in $R$ and $m$.
In particular the right side is uniform in those variables for
bounded $C^1$ endpoint tests and a uniform bound on the partial
derivatives of the phase observations.  All original exact section
events, unvisited constraints, critical endpoints and visited
singular limiting configurations are included with their proved
physical charges.
\end{theorem}
\begin{proof}
On a regular physical word and one true section chamber let
$M=a(x)b(T_R^mx)e^{i\Psi_m(x)}$, and use the label vector $v$
from the proof of Theorem~\ref{thm:av-all-bv}.  Its coefficient is
$wAMv$, where $A=\omega/(4\pi c_*\chi)$ is the unchanged
physical coarea coefficient.  Since $|e^{i\Psi_m}|=1$,
$|M|\le a_0b_0$, and Lemma~\ref{lem:wo-derivatives} gives
\begin{equation}\label{eq:wo-product-derivative}
 |\partial_t M|\le\frac1z
 \left[C_{\rm test,end}(a_1b_0+a_0b_1)
                 +C_{\rm test,phase}\mathcal L_m a_0b_0\right].
\end{equation}
In differentiating $wAM$, the term $M\partial_t(wA)$ is paid
by $a_0b_0$ times the inherited complete base variation.  For
the additional term $wA\partial_tM$ use $0\le w\le1$ and
the full inverse-gradient moment, not a pointwise incidence cutoff:
\begin{align}
 \int_0^{2\pi}\sum_{\text{actual words }v}
       \int_{\mathcal D_v(\theta)}\frac{A_v(t,\theta)}{z}\,dt\,d\theta
 &=\frac1{c_*}\int_{\mathcal C_{m,R}}\frac1{|\nabla_qF|}\,d\nu
 \le\frac1{c_*}\int\frac{I_m}{|\nabla_qF|}\,d\nu
 \le J_{\rm test}.                                  \label{eq:wo-full-moment}
\end{align}
The equality is the one-copy physical coarea formula on the
original cap, first on compact regular pieces and then by
nonnegative exhaustion.  The inequality uses $I_m\ge1$ and
Theorem~\ref{thm:sm-full-multiplier}.  It contains one division
by $c_*$; the separately conditioned inequality from315 is not
inserted and divided by $c_*$ a second time.  In particular all
near-critical positive bulk is retained in this integral.

At a true physical endpoint, the inherited one-sided limits of
every finite-chain collision state exist.  The $C^1$ functions
on $\mathcal X$ are continuous there and bounded, so their
multiplier has a one-sided limit of modulus at most $a_0b_0$.
The complete cap and unvisited-disk charges therefore increase
by at most this factor.  The positive incoming critical sectors
are charged in the same way; no critical value is declared zero.
At an actual visited-grazing boundary the already proved limit
$A\longrightarrow0$ remains zero after bounded multiplication.
This is a trace calculation on each finite chain, not a
deduction from the collision measure of the singular set.

At a regular section cut, $M$ is continuous along the same
physical branch.  The label vector jumps at cost two for an
interior visit and at cost one for an endpoint membership.
Multiplication by $M$ pays at most $a_0b_0$ times those costs.
The coincidence, seam and all-word accounting is precisely the
one-copy accounting in Theorem~\ref{thm:av-all-bv}.  A guard
derivative has already been paid in $D_{\rm base}^w$, and is
not paid a second time in \eqref{eq:wo-full-moment}.

To conclude without assuming a vector-valued integration theorem,
test a finite set of components against independent complex
$C_c^1$ functions of supremum norm at most one.  Integration by
parts on a finite regular chamber exhaustion gives the sum of
the base, section and extra interior charges just computed.
Their full positive total is finite by Lemma~\ref{lem:rbv-endpoints}, Theorem~\ref{thm:rv-exact-bv}, and
\eqref{eq:wo-full-moment}, so the exhaustion passes to the
distributional derivatives.  Maximizing each component test
independently, then increasing the finite label set, gives
\eqref{eq:wo-variation}.  The sum of the $L^1$ norms is at most
the integral of the absolute multiplier against the exact
partition of the original section source.  That integral is
at most $a_0b_0\int w\eta_R(T_R^mx)d\nu_R^*\le a_0b_0$,
which proves \eqref{eq:wo-mass}.
\end{proof}

\begin{corollary}[The weak marked roof operator and its actual norms]
\label{cor:wo-fourier}
For real $\xi\in\mathbb R^2$, $\zeta,\upsilon,s\in\mathbb R$,
define a bilinear form on $C^1(\mathcal X)\times C^1(\mathcal X)$ by
\begin{equation}\label{eq:wo-operator}
 \begin{split}
 \langle\mathcal K_{m,R}^{w;\xi,\zeta,\upsilon}(s)a,b\rangle
 =\int &w\eta_R(T_R^mx)a(x)b(T_R^mx)\\
 &\cdot e^{i(\xi\cdot K^{\rm c}_{m,R}(x)+
       \zeta A_{m,R}(x)+\upsilon m+sL_{m,R}(x))}
          \,d\nu_R^*(x).
 \end{split}
\end{equation}
It defines a bounded linear map $C^1(\mathcal X)\to C^1(\mathcal X)^*$,
where the dual acts linearly on the second test, and for $s\ne0$,
\begin{equation}\label{eq:wo-weak-bound}
 \|\mathcal K_{m,R}^{w;\xi,\zeta,\upsilon}(s)\|_{C^1\to(C^1)^*}
 \le\min\left\{1,
       \frac{D_{\rm all}^w+2J_{\rm test}C_{\rm test,end}}{|s|}\right\}.
\end{equation}
The bound is uniform in $m,R$ and all displayed real mark
frequencies.  Insertion of $e^{i\Psi_m}$ adds at most
$J_{\rm test}C_{\rm test,phase}\mathcal L_m$ to the numerator.
For any bounded label coefficients, the corresponding tested
roof density has translation error at most
$\mathcal V_m^w(a,b,\Psi)|h|$ in $L^1$.
\end{corollary}
\begin{proof}
Sum the densities in \eqref{eq:wo-measures} with coefficients
$e^{i(\xi\cdot k+\zeta n+\upsilon m)}$.  The absolute $L^1$
and variation sums in Theorem~\ref{thm:wo-bv} justify this
operation and bound the resulting density and derivative.
For a compactly supported BV function $f$ and $s\ne0$,
distributional integration by parts gives
$|\int e^{ist}f(t)dt|\le\|Df\|_{\rm TV}/|s|$.
The $L^1$ bound gives the other part of the minimum.
If both $C^1$ test norms are at most one, then
$a_0b_0\le1$ and $a_1b_0+a_0b_1\le2$; the displayed
constant follows without discarding either partial-derivative
norm.  The phase addition follows from the same calculation.
The translation assertion follows by integrating the derivative
measure along a roof displacement and taking its total variation.
\end{proof}

\paragraph{Relation to the original local problem.}
At fixed $m$ the phase $\upsilon m$ records the collision mark;
it does not create an extra coordinate beyond the original
four-coordinate slice $(k_1,k_2,n,t)$.  Formula
\eqref{eq:wo-operator} is evaluated on the actual physical
source with the original entrance law and terminal membership.
The full chain phase is an optional observation with the exact
norm stated above, and no derivative of an arbitrary Borel
trajectory selector is claimed.  Such selectors are covered
by the positive likelihood estimate of
Theorem~\ref{thm:av-likelihood} instead.

Estimate \eqref{eq:wo-weak-bound} supplies a weak tested
$1/|s|$ roof estimate.  It is neither a strong anisotropic
transfer-operator estimate nor a decay estimate in the number
of returns.  In particular $1/|s|$ is not integrable at infinity;
absolute Fourier inversion on an expanding frequency range
does not follow from this bound alone.  The arithmetic peak
expansion, uniform spectral remainder, microscopic pointwise
local theorem and long-time four-coordinate inverse require
their stated additional inputs.  No supremum or ordered limit
in the original target is exchanged here.  The endpoint
stationary-action derivative used in the proof does not
replace the actual long-time forward-orbit Jacobian in that
inverse problem.
